% Lernkarten -- MODULNAME (Modul NUMMER), Lektion N
% Quelle: eigene handschriftliche Notizen (z. B. OneNote), NICHT das Glossar.
%
% Format je Karte:  \lernkarte{Vorderseite}{Rueckseite}
%   Vorderseite = roter Titel/Begriff aus den Notizen (Ueberschrift)
%   Rueckseite  = weisser Inhalt darunter (Definition/Formel/Regel)
%
% Diese Datei wird NICHT als eigenstaendiges PDF kompiliert. generate_docs.py
% (extract_lernkarten) parst nur die \lernkarte-Makros; Vorder- und Rueckseite
% werden einzeln in Flip-Karten (Web) und Anki-Karten gerendert. Der Inhalt
% beider Argumente ist reines LaTeX -- Mathe inline mit $...$ oder abgesetzt
% mit \[...\]; Aufzaehlungen mit enumitem; Tabellen mit array/booktabs.
%
% Verfuegbare Pakete beim Rendern: amsmath, amssymb, enumitem, array, booktabs.
% Kommentarzeilen (mit % am Zeilenanfang) werden ignoriert und koennen daher
% gefahrlos das Wort \lernkarte enthalten.
%
% \lernabschnitt{1.1 Reelle Zahlen}  setzt den aktuellen Abschnitt; er wird
% klein ueber den Titel jeder folgenden Karte gesetzt.
%
% \lernkarte[id]{Titel}{Rueckseite}  Die optionale [id] ist ein STABILER,
% inhaltsunabhaengiger Bezeichner; daraus wird die Anki-GUID gebildet. Schema:
% <modul>-<lektion>-<laufnummer>, z. B. [61211-1-01]. Titel/Rueckseite duerfen
% sich spaeter aendern, ohne den Anki-Lernfortschritt zu verlieren -- die id
% NICHT aendern und je Karte eindeutig vergeben (neue Karte = neue id).

\lernabschnitt{1.1 Titel des Abschnitts}

\lernkarte[MODUL-N-01]{Titel der ersten Karte}{%
  Inhalt der Rueckseite -- z. B. eine Definition mit Formel:
  \[ a^2 + b^2 = c^2 \]
}

\lernkarte[MODUL-N-02]{Titel der zweiten Karte}{%
  \begin{itemize}[topsep=2pt,itemsep=1pt]
    \item erste Eigenschaft
    \item zweite Eigenschaft
  \end{itemize}
}
